\documentclass[11pt,a4paper]{article}
\usepackage[margin=25mm]{geometry}
\usepackage{amsmath,amssymb,amsthm,hyperref}
\newtheorem{theorem}{Theorem}
\newtheorem{lemma}{Lemma}
\newcommand{\gp}[1]{\langle #1\rangle}
\title{Correlation reductions and the exact square-torus parameters}
\author{Human Compression project}\date{}
\begin{document}\maketitle
\begin{abstract}
We supply the definite diagram-(4) reduction stated on physical page 9
of Fernandes--Leemans--Weiss, and the definite presentation statements
accompanying the experiments concerning diagram (5). Explicit inverse
presentation maps give the parameter transformations. For every existing regular fork hypertope, its correlation
extension has the intersection property and is the stated universal
polytope. The group identities do not assert existence of a hypertope for
every period or degenerate quotient.
\end{abstract}

\section{The literal groups and their correlation}
Write $a=\rho _0$, $b=\rho _1$, $c=\rho _2$, $d=\rho _3$, with $b$ the
central node of the fork. Its Coxeter relations are
\begin{equation}\label{fork}
 a^2=b^2=c^2=d^2=(ab)^4=(bc)^4=(bd)^3
 =(ac)^2=(ad)^2=(cd)^2=1.
\end{equation}
Let $H_s^0$ be this presented group with $(abcb)^s=1$, and let $H_s^1$
have instead $(abc)^{2s}=1$, where $s$ is a positive integer. These are
the source's axial $(s,0)$ and diagonal $(s,s)$ square-torus residue
relations, respectively. The assertion here concerns these entire
presentations, not a finite image of them.

The substitution $\alpha:a\leftrightarrow c$, fixing $b,d$, preserves
\eqref{fork}. Put $R=abcb$ and $Q=abc$. Then
\[
 \alpha(R)=cbab=bR^{-1}b,\qquad \alpha(Q)=cba=Q^{-1}.
\]
It therefore preserves each of the two additional normal closures.
It induces an automorphism of $H_s^\epsilon$ with square one, whether
or not that presented group already defines a nondegenerate geometry.
Form the external semidirect product
\[
 E_s^\epsilon=H_s^\epsilon\rtimes_\alpha\gp{t},\qquad t^2=1.
\]
Thus $tat=c$, $tbt=b$, $tdt=d$. The external factor has order two even
in a collapsed small-parameter group. All statements about actual
hypertope types below impose their stated nondegeneracy hypotheses.

\begin{theorem}[Exact universal presentation reduction]\label{universal}
On the ordered generators $(t,a,b,d)$, $E_s^0$ is the Coxeter group
$[4,4,3]$ with the extra relation $(tab)^{2s}=1$, and $E_s^1$ is
$[4,4,3]$ with the extra relation $(taba)^{2s}=1$. In particular their
marked local square-torus presentations are respectively $(s,s)$ and
$(2s,0)$. In both cases $H_s^\epsilon$ is exactly the kernel of the
character which sends $t$ to $-1$ and $a,b,d$ to $1$.
\end{theorem}
\begin{proof}
Eliminate $c$ using $c=tat$. The relation $(ac)^2=1$ becomes
$(at)^4=1$. The relations involving $bc$ and $cd$ are conjugates by $t$
of those involving $ba$ and $ad$; $t$ commutes with $b,d$. The remaining
relations are exactly the string $[4,4,3]$ relations on $(t,a,b,d)$.
For the periods the following are word identities, not order tests:
\begin{align}
 abcb&=abtatb=(abt)^2=t(tab)^2t,\label{axial}\\
 (abc)^2&=abacbc=(abat)^2=(taba)^{-2}.\label{diagonal}
\end{align}
In the first line commuting $t$ and $b$ gives $(abt)^2$.
In the second line commute $a,c$, then use
$tat=c$ and $tbt=b$. Taking the indicated powers proves equivalence of
the extra relators and their full normal closures.

Conversely start with either stated string presentation, set $c=tat$,
and retain $t$. The preceding steps give every fork relation, the correct
period, and the conjugation rules for $\alpha$. The resulting homomorphisms
between that presentation and $E_s^\epsilon$ are inverse on all generators.
This proves isomorphism of the presented groups, including the universal
upper direction. Every defining string relator has even $t$-parity, so
the character exists and is onto. In the semidirect product its kernel
is precisely $H_s^\epsilon$. Equivalently, the Schreier transversal
$\{1,t\}$ gives generators $a,c,b,d$ and the already displayed fork
presentation; no omitted conjugate relator is required.
\end{proof}

The usual square-torus convention for a string marking $(x,y,z)$ is
$(xyz y)^q=1$ for $(q,0)$ and $(xyz)^{2q}=1$ for $(q,q)$.
Thus \eqref{axial} and \eqref{diagonal}, with $(x,y,z)=(t,a,b)$,
give exactly the two parameter transformations in the theorem. This
uses the same torus convention as the source and the retained square
theory; it does not identify a torus by its group order. The convention
can also be checked on the square plane: take
\[
 a(x,y)=(-x,y),\quad c(x,y)=(x,-y),\quad
 b(x,y)=(1-y,1-x),\quad t(x,y)=(y,x).
\]
Here $abcb$ translates by $(-2,0)$, $(abc)^2$ by $(-2,2)$, and
$taba$ by $(1,-1)$. Rotation-conjugate period vectors give the
complete two-dimensional period lattices. The string grid has basis
$(1,-1),(1,1)$, so its diagonal $(s,s)$ lattice is the fork's axial
$2s\mathbb Z^2$, while its axial $(2s,0)$ lattice is generated by
$(2s,-2s),(2s,2s)$. This confirms both the scaling and the exchange
of axial and diagonal labels. These coordinates are only the local
square-plane model, not a purported model of the rank-four group.

\section{Intersection property and the geometric assertion}
A set of involutions has the intersection property if its generated
subgroups satisfy $H_I\cap H_J=H_{I\cap J}$ for all subsets of its
marking. The following elementary form of the usual parabolic patching
argument suffices; no geometry is inferred just from an abstract group
isomorphism. We retain the standard equivalence between string C-groups
and regular abstract polytopes, and the regular-hypertope
maximal-parabolic coset representation, as explicit foundational inputs.

\begin{lemma}[Parabolic patching]\label{patch}
If every maximal proper generated parabolic has the intersection property,
and every two distinct such parabolics intersect in their common generated
subgroup, then the entire marking has the intersection property.
\end{lemma}
\begin{proof}
Let $G_I,G_J$ be proper generated subgroups. Choose a missing marker $i$
of $I$ and a missing marker $j$ of $J$. If $i=j$, the intersection property
in that maximal parabolic applies. Otherwise an element of their
intersection lies in the common parabolic $G_T$, where $T$ omits $i,j$.
Applying the intersection property first inside $G_{\bar i}$ and then
inside $G_{\bar j}$ puts it in $G_{I\cap T}$ and $G_{J\cap T}$.
Both are in $G_T$, which inherits the intersection property from a maximal
parabolic. Their intersection is $G_{I\cap J}$, as required. Cases with
$I$ or $J$ the full marking are immediate.
\end{proof}

\begin{theorem}[From each existing fork to its polytope]\label{geometry}
Suppose $(H_s^\epsilon;a,b,c,d)$ is the actual marked group of a regular
hypertope of source diagram (4), with the prescribed square-torus residue
and the prescribed cubical residues. Then $(E_s^\epsilon;t,a,b,d)$ is
a string C-group. Its polytope is respectively
\[
 \bigl\{\{4,4\}_{(s,s)},\{4,3\}\bigr\},\qquad
 \bigl\{\{4,4\}_{(2s,0)},\{4,3\}\bigr\},
\]
and is universal with those marked local defining presentations. The
given fork hypertope is recovered from its parity kernel and marking.
\end{theorem}
\begin{proof}
The marked group of the given regular hypertope has the intersection
property. In $E=E_s^\epsilon$ put $L=\gp{t,a,b}$ and $K=\gp{a,b,d}$.
The semidirect normal form gives
\[
 L\cap H_s^\epsilon=\gp{a,b,c},\qquad
 \gp{t,a}\cap H_s^\epsilon=\gp{a,c}.
\]
Indeed the listed kernel generators and $t$ generate the respective
subgroups, and $t$ normalizes each listed kernel. Consequently
\[
 \gp{t,a}\cap\gp{a,b}
 =\gp{a,c}\cap\gp{a,b}=\gp a.
\]
The two involutions $a,c$ are independent commuting markers, so $ta$
has order four. The remaining maximal subgroup of $(t,a,b)$ is
$\gp{t,b}=\gp t\times\gp b$. Its intersections with $\gp{t,a}$ and
$\gp{a,b}$ are respectively $\gp t$ and $\gp b$, using parity and
$\gp{a,c}\cap\gp b=1$. Each two-marker subgroup has the intersection
property, so Lemma~\ref{patch} makes $(t,a,b)$ a C-group.
The other end deletion $(a,b,d)$ inherits the intersection
property and its string diagram directly from $H_s^\epsilon$. Finally
\[
 L\cap K=\gp{a,b,c}\cap\gp{a,b,d}=\gp{a,b}.
\]
For completeness the other two maximal parabolics are
$M_1=\gp{t,b,d}=\gp t\times\gp{b,d}$ and
$M_2=\gp{t,a,d}=\gp{t,a}\times\gp d$. The first product is direct by
parity; the second uses $\gp{a,c}\cap\gp d=1$ and the displayed
commutations. Thus they too have the intersection property. In the
ordering $(M_0,M_1,M_2,M_3)=(K,M_1,M_2,L)$ the six intersections are
\[
\begin{array}{c|cccccc}
(i,j)&(0,1)&(0,2)&(0,3)&(1,2)&(1,3)&(2,3)\\ \hline
M_i\cap M_j&\gp{b,d}&\gp{a,d}&\gp{a,b}&\gp{t,d}&\gp{t,b}&\gp{t,a}.
\end{array}
\]
To verify the last three, both groups contain $t$, so intersect their
parity kernels $\gp{b,d}$, $\gp{a,c,d}$, $\gp{a,b,c}$ inside $H$ and
then reattach $\gp t$. The first three follow in the same way with
$K\subset H$. The original intersection property supplies every kernel
intersection. Lemma~\ref{patch} proves the full string intersection
property. Its cubical vertex figure is the given $(a,b,d)$ residue.
Its facet group is the correlation extension of the given entire
$(a,b,c)$ torus group. Applying the rank-three presentation calculation
of Theorem~\ref{universal} identifies the facet with the stated torus,
without an additional local quotient or period collapse.

Theorem~\ref{universal} supplies exactly the global string relations and
the defining facet period, with no further relators. Hence every regular
polytope of those marked local types is a quotient, and the displayed
polytope is universal. Conversely its exact parity kernel is the original
marked group. A regular hypertope with that marking is its standard
maximal-parabolic coset geometry, so this recovers the original geometry,
not just an unmarked group order.
\end{proof}

The type exchange $a\leftrightarrow c$ is an actual correlation of the
given regular fork geometry: on its maximal-parabolic cosets apply
$\alpha$ and exchange the corresponding two types. This preserves
coset intersection and is an involution. The new rank-four generating
list is $t,a,b,d$. The source's three-letter parenthesis omits $\rho _3$;
$t,a,b$ generates the facet extension, while the complete rank-four
assertion requires $d=\rho _3$ as well.

\section{Diagram (5): definite algebra, separate geometric experiment}
For diagram (5), keep the labels $a,b,c$ of its square-torus path, with
$b$ central, and put $d$ at the other vertex of the diamond. Its Coxeter
relations are
\begin{equation}\label{diamond}
 a^2=b^2=c^2=d^2=(ab)^4=(bc)^4=(ad)^3=(cd)^3
 =(ac)^2=(bd)^2=1.
\end{equation}
This agrees with the actual source drawing: the two edges at $b$ have
label four and the two at $d$ have the undisplayed label three.
Let $D_s^0,D_s^1$ add, respectively, $R^s=1$ and $Q^{2s}=1$ to
\eqref{diamond}. The same substitution $a\leftrightarrow c$, fixing
$b,d$, is an involutory automorphism of both groups.

\begin{theorem}[Diamond-to-fork presentations]\label{diamondfold}
The extension $D_s^0\rtimes\gp t$, marked $(t,a,b,d)$ with $a$ central,
has exactly the diagram-(4) fork presentation with torus parameter $(s,s)$.
The extension $D_s^1\rtimes\gp t$ has exactly that fork presentation with
torus parameter $(2s,0)$. In both cases $D_s^\epsilon$ is its exact
$t$-parity kernel. These statements hold for all positive integers $s$
as isomorphisms of the displayed presented groups.
\end{theorem}
\begin{proof}
Again eliminate $c=tat$. The relation $(ac)^2$ gives $(ta)^4$; $(cd)^3$
is the $t$-conjugate of $(ad)^3$; $(bc)^4$ is the conjugate of $(ba)^4$.
The remaining Coxeter relations give the fork with centre $a$, arms
$t,b$ of label four and arm $d$ of label three. The period words transform
by exactly \eqref{axial}--\eqref{diagonal}. Thus they become $(tab)^{2s}$
and $(taba)^{2s}$, respectively. Now $t,a,b$ is the fork's toroidal path,
so these are its diagonal $(s,s)$ and axial $(2s,0)$ periods. The inverse
substitution gives all original relations and the conjugation rules,
proving both directions and the exact kernel assertion. In particular the
source's conjugated word is correct: $\alpha(R)=cbab=(cbt)^2$ and
$cbt=tab$. Conjugating $R^s=1$ therefore gives precisely
$(cbt)^{2s}=1$.
\end{proof}

If the original diamond marking is a C-group, the extended fork marking
also has the intersection property. The six maximal-parabolic
intersections in the proof of Theorem~\ref{geometry} are unchanged.
The one modification is $M_2=\gp{t,a,d}$: it is now a rank-three string
marking of type $\{4,3\}$, rather than a direct product. Its three
two-marker intersections follow from
$\gp{t,a}\cap H=\gp{a,c}$, the original intersection property and
$[t,d]=1$, exactly as for $\gp{t,a,b}$. Lemma~\ref{patch} applies first
there and then to the full marking. This still does not prove flag
transitivity of the new nonstring coset incidence system. The source
explicitly calls hypertope production through this correlation experimental;
the definite group and period assertions above do not remove that
geometric qualification.

\section{Exact scope and converse limitations}
The group isomorphisms hold for every positive integer $s$. The geometric
determination is quantified over each existing regular universal fork
hypertope with its stated actual residues. It determines its marked group
and coset geometry from the corresponding universal polytope. It does
not claim that every period in the displayed presentations gives a
hypertope, or that every quotient preserves its local types, intersection
property or flag transitivity. In a proposed converse construction those
properties must be established, not deduced merely from the parity map.
For example, the separately proved bipartite halving theorem applies to a
regular lattice polytope with a bipartite vertex graph; its actual lattice
and incidence hypotheses must be checked before using it. No such
unchecked hypothesis is inserted into Theorem~\ref{geometry}.

This repairs the earlier combined documentary disposition for diagrams
(4) and (5). Diagram (4) has the definite universal reduction just proved.
Diagram (5)'s exact stated presentation and parameter transformations
are proved separately in Theorem~\ref{diamondfold}; its experimental
hypertope-existence wording remains explicitly qualified. Neither an
unproved flag-transitivity claim nor the order-768 future-work example is
silently added to this supplement.

The source is M. E. Fernandes, D. Leemans and A. Ivi\'c Weiss,
\emph{Hexagonal extensions of toroidal maps and hypermaps}, the exact
retained 21-page 2018 edition, physical pages 9--10. The correlation idea
and parameter reduction are theirs. The work here writes the inverse
presentation maps and the missing marked-intersection argument, reusing
the established Schreier and string-C-group principles. It claims no
new original admission, historical priority, question solution or Lean
formalization. No presentation enumeration or previously accepted
geometric census is repeated.
\end{document}
